\section{A Wedding, a Garment, and an Evening Prayer}
\label{sec:opening}

The Nineteenth Sunday after Pentecost, \latin{Dominica decima nona post
Pentecosten}, is a Sunday of the second class in the Proper of Time of the 1962
\work{Missale Romanum}. Its Mass begins in the lower left-hand column of
p.~411 of the typical edition, directly beneath the Postcommunion of the
Eighteenth Sunday, and ends near the head of p.~412, where the Twentieth
Sunday begins; its ten elements carry the marginal numbers 1679--1688. In
2026 it falls on 4~October, the first Sunday of the month. The calendar's
feast of that day, St Francis, is a feast of the third class and gives way
entirely, since a Sunday of the second class admits no commemoration except
that of a feast of the second class (\work{Rubricae generales} 111~b;
\work{Rubricae generales Missalis} 434~b). The Mass therefore has one Collect,
one Secret and one Postcommunion. It is said in green (RG 127~b), with Gloria,
with the Credo that the rubrics require on every Sunday (RGMR 475~a), and with the Preface of the Most Holy Trinity, which
the Missal prints as a direction after the Secret and which the general rubrics
appoint for the Sundays of the second class outside the Christmas and Easter
seasons (RGMR 494~b).

The first Sunday of October also carries by law the external solemnity of the
Most Holy Rosary (RGMR 358~b). A church that keeps it may say one sung and one
low Mass, or two low Masses, of the Rosary as votive Masses of the second class
(RGMR 360), and in those Masses this Sunday's Collect, Secret and Postcommunion
are said as a commemoration. Every other Mass of the day is the Mass studied
here.

That Mass opens with God speaking in the first person. \latin{Salus populi ego
sum, dicit Dominus}; in the English of the 1861 Philadelphia hand missal, ``I am
the Saviour of my people, saith the Lord: in whatever distress they call on me,
I will hear them: and will be their Lord for ever.'' The psalm verse that
follows is the opening of a long psalm in which Israel is told its own history:
``Attend, O my people, to my law.'' The Collect asks that everything adverse,
\latin{universa nobis adversantia}, be shut out, so that those who pray it,
\latin{mente et corpore pariter expediti}, may carry out \latin{quae tua sunt},
the things that are God's. The
Epistle, from the fourth chapter of Ephesians, commands the renewal of the mind
and the putting on of ``the new man, who according to God is created in
justice and holiness of truth'', and then names what that looks like in four
plain acts: truth to the neighbour, anger ended before the sun goes down, no
foothold for the devil, and work that has something to give. The Gradual is an
evening prayer: ``Let my prayer be directed as incense in thy sight; the
lifting up of my hands, as evening sacrifice.'' The Alleluia calls for praise,
for invocation, and for the Lord's deeds to be declared among the Gentiles.

The Gospel is the parable of the king who made a marriage for his son (Mt
22:1--14). The invited guests will not come; invited a second time, they go
off to farm and merchandise, and the rest seize the servants and kill them. The
king sends his armies and burns their city, and his servants go out into the
highways and bring in everyone they find, ``both bad and good''. Then the king
comes in to see the guests, and finds one man without a wedding garment. He
calls him ``Friend'', the man has nothing to say, and he is bound and cast out:
``For many are called, but few are chosen.'' The Offertory takes up the
Introit's word \latin{tribulatio}: ``If I shall walk in the midst of
tribulation, thou wilt quicken me.'' The Secret offers the gifts \latin{oculis
tuae maiestatis}, before the eyes of God's majesty, and asks that they be
\latin{salutaria}, saving. The Communion sets a command beside a wish, ``Thou
hast commanded thy commandments to be kept most diligently. O! that my ways may
be directed to keep thy justifications'', and the Postcommunion asks two things
of God's \latin{medicinalis operatio}, his healing work: that it free us from
our \latin{perversitates}, and that it make us \latin{inhaerere mandatis}, cleave
to his commandments.

The Gospel ends with a man who has come in and is found wanting, and the rest
of the formulary is full of calling, hearing and keeping. Three questions arise
from it, and each gives one of the readings that follow. The first asks what
the guest who has come in must be wearing. It is heard from the end of the
Gospel and from the Epistle's new man, and St Gregory the Great, St Augustine
and St Jerome carry it: the garment is charity, the new man put on, Christ
himself, shown in the Epistle's plain acts and asked for in the Communion's
wish. The second asks who was called, by whom, and what became of the call
when it was refused. It is heard from the beginning of the Gospel, from the
Introit's psalm of Israel's history and from the Alleluia's command to the
nations, and St John Chrysostom, St Irenaeus and St Hilary of Poitiers carry
it: the call passed from the first invited to the highways. The third asks
where the feast is now and how the assembly at this Mass takes its place at
it. It is heard from the Gospel's crowded hall and from the chants and prayers
of the Mass itself, the evening prayer of the Gradual, the Offertory's walk
through tribulation, the gifts of the Secret and the Communion's wish. St
Augustine and St Gregory carry it, with St John Chrysostom's address to those
who have partaken of the mysteries: the feast is the Church now, at the Lord's
table, where the King comes in to see his guests. Bl.\ Ildefonso Schuster,
commenting on this Mass, reads its Gradual as the evening oblation of the
Christian life and its Secret and Postcommunion of the Eucharist. Every element of the formulary has a place in each reading,
and each closes with four senses of its own, literal, allegorical, moral and
anagogical. The three readings answer different questions about the same texts
and do not contradict one another, though the witnesses within each do not
always agree.

\subsection{The appointed elements at a glance}

\begin{maptable}
Introit & \latin{Salus populi ego sum}; verse Ps 77:1 (78:1) &
God speaks: he is the salvation of the people, hears them from every
tribulation, and will be their Lord for ever. The verse opens Asaph's psalm of
Israel's history: ``Attend, O my people, to my law.''\\
Collect & \latin{Omnipotens et misericors Deus} &
In the 1861 English, ``favourably defend us from all adversity: that being
free both in soul and body, we may with security of mind perform thy
service''. Its word \latin{expediti}, set free, returns in the
Postcommunion.\\
Epistle & Eph 4:23--28 &
Renewal in the spirit of the mind; the new man ``created in justice and
holiness of truth''; truth to the neighbour, ``for we are members one of
another''; anger ended before sunset; no place for the devil; labour ``that he
may have something to give to him that suffereth need''.\\
Gradual & \latin{Dirigatur oratio mea}; Ps 140:2 (141:2) &
``Let my prayer be directed as incense in thy sight; the lifting up of my
hands, as evening sacrifice.''\\
Alleluia & \latin{Confitemini Domino}; Ps 104:1 (105:1) &
``Give glory to the Lord, and call upon his name: declare his deeds among the
Gentiles.''\\
Gospel & Mt 22:1--14 &
The king's marriage feast for his son: the invitation twice refused, the
servants killed, the city burned, the highways, the hall filled with bad and
good, and the guest without a wedding garment. ``Many are called, but few are
chosen.''\\
Offertory & \latin{Si ambulavero}; Ps 137:7 (138:7) &
The psalmist walks in the midst of tribulation and trusts God to keep him
alive; in the Missal's form God's hand and right hand are still to act.\\
Secret & \latin{Haec munera, quaesumus, Domine} &
In the 1861 English, ``that the offerings we bring before thy divine Majesty
may avail to our salvation''. The Latin offers them \latin{oculis tuae
maiestatis}, to the eyes of God's majesty.\\
Communion & \latin{Tu mandasti}; Ps 118:4--5 (119:4--5) &
``Thou hast commanded thy commandments to be kept most diligently. O! that my
ways may be directed to keep thy justifications.''\\
Postcommunion & \latin{Tua nos, Domine, medicinalis operatio} &
In the 1861 English, ``May the healing efficacy of these thy mysteries, O
Lord, mercifully free us from our perverseness, and make us always obedient to
thy commandments.''\\
\end{maptable}
